\documentclass{article}
\usepackage{amssymb, amsmath}
\usepackage[left=1.5cm, right=2cm]{geometry}
\begin{document}

\noindent Gianpaolo Girasoli\\
CSE3400 Hw2Part2

\textbf{Problem 1}\begin{enumerate}
	\item This MAC is secure since it uses the PRF $F$. The outer PRF returns a deterministic, pseudorandom string and is therefore also a secure MAC. This means that the attacker's probability of generating a forgery is $\frac{1}{2^n}$, resulting in a negligible advantage.
	\item This scheme is secure. Since $m$ is appended with different bit strings ($1^{n/2}$ and $0^{n/2}$ for the first and second halves respectively), the secure pseudorandom strings generated by each $F_k$ cannot be interchanged. Namely, $m||1^{n/2}$ and $m||0^{n/2}$ will always be different. Even if an attacker tries to rearrange computed tags from the oracle, since the appended bits are strictly placed, there is no way the attacker can create a forgery. Thus, the probability of generating one is $\frac{1}{2^n}$, resulting in a negligible advantage.
	\item This scheme is insecure. For an attack, an attacker can query the MAC oracle for a tag over a message of $n-2$ bits of 1s and append 01. This results in a tag $t$, with $F_k(\overline{m})$ being computed. This is also a correct tag over a message of $n-2$ bits of 0s, appended with 10 (resulting in $F_k(m)$ being evaluated, resulting in a tag that equals $t$). Therefore, the probability of making a correct forgery is 1, and the advantage of the attacker is $1-\frac{1}{2^n}$, which is non-negligible. Therefore, the probability of making a correct forgery is 1, and the advantage of the attacker is $1-\frac{1}{2^n}$, which is non-negligible.
\end{enumerate}

\textbf{Problem 2}\begin{enumerate}
	\item $h'$ is not CRHF. Since $h_2$ is not CRHF (even though it is OWF), an attacker can find a collision via messages $m$ and $m'$, where $m=m_1||m_2$ and $m'=m_1||m_3$ such that $h_2(m_2)$ and $h_2(m_3)$ collide (since it is not guaranteed that it is hard to find a collision for a OWF, and $m_2$ and $m_3$ are not random, being chosen by the attacker). Since an attacker can find a collision for $h'$, the success probability is 1 and the attacker's advantage is $1-\frac{1}{2^n}$, which is non-negligible.
	\item $h''$ is OWF, since the outermost layer is $h_2$, which is also OWF. In the OWF game, an attacker does not get the outputs for $h_1$ (which constitutes the input for $h_2$) and therefore cannot find a collision, resulting in a success probability of $\frac{1}{2^n}$, which is negiligible.\\
		However, $h''$ is not CRHF. Since the outermost layer is OWF but not CRHF, it is not hard for an attacker to find a collision. Since they are playing the CRHF, they can just exhaustive search messages until they find a collision for $h_2$. Therefore, the success probability is $1-\frac{1}{2^n}$, which is non-negligible.
	\item Sponge's proposal will preserve the security. Since $h_1$ is already CRHF and is the outermost layer, it will preserve being CRHF which is required for the application. Alice's proposal, using $h_2$ as the outermost layer, is not CRHF. An attacker, in the CRHF game, can simply brute force inputs to $h_2$ and find a collision since it is not CRHF. This results in a success probability of 1, with a non-negligible advantage of $1-\frac{1}{2^n}$.
	\item\begin{itemize}
		\item[(a)] The number of nodes at the leaf level is 32 (one node for each file and its hash). The total size is 47 (32 leaf nodes, 5 layers of hashes including the leaf nodes and the root node).
		\item[(b)] The proof of inclusion is $h_{1-8},h_{9-12},h_{14}$, and $h_{15-16}$. The size of this is 5. Eve can verify this PoI by computing $h(h(f_{13})||h_{14})$ to get $h_{13-14}$, then hash that with $h_{15-16}$ to get $h_{13-16}$. Then, hash that with $h_{9-12}$ to get $h_{9-16}$, and then finally checking if that hashed with $h_{1-8}$ is the same as the root node of the Merkle Tree.
	\end{itemize}
\end{enumerate}

\end{document}
